\chapter{Curricula}
\label{chp:INT}



\section{Österreichische Sanitäter-Ausbildungsverordnung}

\HeadingSub{SanAV}

\subsection{Ausbildungsziele des Modul 1}


\begin{WidePar}
% \scriptsize

\begin{multicols}{2}

Ausbildungsziele für österreichische Rettungssanitäter gem. der \emph{Verordnung
der Bundesministerin für Gesundheit und Frauen über die Ausbildung zum
Sanitäter - Sanitäter-Ausbildungsverordnung - San-AV StF: BGBl. II Nr. 420/2003
-- Anlage 1, Modul I}.

\begin{enumerate}
  
  \item \TabHeading{Erste Hilfe und erweiterte Erste Hilfe}
  
  (Bestandteil des Sachgebiets 'Sanitätshilfe')
  
  \begin{inparadesc}
    \item[Stundenanzahl:] 14
    \item[Lehrkraft:] Notarzt, Arzt für Allgemeinmedizin, approbierter Arzt,
    Facharzt, Lehrsanitäter, fachkompetente Person
  \end{inparadesc}
  \begin{enumerate}
    \item
    Notwendigkeit, Verpflichtung der Hilfeleistung
    \\
    \noindent\FullRef{topic:behandlungspflichtJUS},
    \FullRef{topic:behandlungspflicht-ehi}
    \item
    Unfallverhütung
    \\
    \FullRef{topic:verhalten-notfallort},
    \FullRef{topic:gefahrenzone}
    \item
    Rettungskette
    \\
    \noindent\FullRef{topic:rettungskette}
    \item
    Gefahrenzonen
    \\
    \noindent\FullRef{topic:gefahrenzone}
    \item
    Bergen (Wegziehen, Rautekgriff aus dem  Auto, Sturzhelm)
    \\
    \FullRef{topic:wegziehen},
    \FullRef{topic:rautekgriff}, 
    \FullRef{topic:sturzhelmabnahme}
    \item
    Notfallpatient, Notfalldiagnose, Notfallhilfe
    \\
    \FullRef{topic:notfalldiagnose-bewusstlosigkeit},
    \FullRef{topic:notfalldiagnose-atem-kreislaufstillstand},
    \FullRef{topic:starke-blutung},
    \FullRef{topic:schock-ehi}
    \item    Kontrolle der Lebensfunktionen
    \\
    \noindent\FullRef{topic:lebensfunktionen-kontrolle}
    \item    Notfalldiagnose Bewusstlosigkeit
    \\
    \FullRef{topic:notfalldiagnose-bewusstlosigkeit}
    \item    Notfalldiagnose  Atem-Kreislaufstillstand
    \\
    \noindent\FullRef{topic:notfalldiagnose-atem-kreislaufstillstand}
    \item    Starke Blutung (Blutstillung, Fingerdruck,  Druckverband, Abbindung)
    \\
    \noindent\FullRef{topic:starke-blutung}
    \item    Schock (Schockbekämpfung,  Lagerungen)
    \\
    \noindent\FullRef{topic:schock-ehi}
    \item    Wunden und Wundverbände (Tierbisse, Insektenstiche, Verätzungen,
    Verbrennungen, Erfrierungen, Unterkühlung) 
    \\
    Wunden: \FullRef{topic:wunden-ehi} \\
    Wundverbände: \FullRef{topic:wundverbaende}
    \item    Quetschungen
    \\
    \FullRef{topic:wunden-einteilung}
    \item    Gelenksverletzungen
    \\
    Verstauchung: \FullRef{topic:verstauchung-eh},
    \\
    Verrenkung: \FullRef{topic:verrenkung-eh}
    \item
    Knochenbrüche, allgemein
    \\
    \noindent\FullRef{topic:knochenbruch-eh}
    \item
    Brustkorbverletzungen
    \\
    \FullRef{topic:thoraxtrauma}
    \item    Plötzlich auftretende Erkrankungen
    \\
    \FullRef{chp:MED}
    \item    Vergiftungen
    \\
    \noindent\FullRef{topic:vergiftungen-ehi}
    \item    Transport
    \\
    \FullRef{topic:rettungskette}
  \end{enumerate}
  \item \TabHeading{Hygiene}
  
  (Bestandteil des Sachgebiets 'Sanitätshilfe')
  
  \begin{inparadesc}
    \item[Stundenanzahl:] 2
    \item[Lehrkraft:] Arzt für Allgemeinmedizin, approbierter Arzt,
    Facharzt, Angehörige des gehobenen Dienstes für Gesundheits- und Krankenpflege,
    Hygienefachkraft,  Lehrsanitäter
  \end{inparadesc}
  
  \begin{enumerate}
    \item    Persönliche Hygiene
    \\
    \FullRef{topic:persoenliche-hygiene}
    \item    Grundbegriffe der Desinfektion
    \\
    \FullRef{topic:desinfektion-grundzuege}
    \item    Grundbegriffe der  Sterilisation
    \\
    \FullRef{topic:sterilisation-grundzuege}
    \item    Entsorgung von infektiösem Abfall
    \\
    \FullRef{topic:entsorgung-infektioeser-abfall}
    \item    Allgemeine Infektionslehre
    \\
    Erreger: \FullRef{topic:infektion-erreger},
    
    Übertragungswege: \FullRef{topic:infektion-uebertragungswege},
    
    Kreuzinfektion: \FullRef{topic:kreuzinfektion}.
    \item
    Vorgehen bei Verletzungen des Personals
    \\
    Nadelstichverletzungen: \FullRef{topic:nadelstichverletzung}
    \item    Hygienemaßnahmenplan
    \\
    \FullRef{topic:hygieneplan}
    \item    Infektionstransport
    \\
    \Abk{MRSA} und andere resistente Keime: \prettyref{m:mrsa-transport}
  \end{enumerate}
  \item \TabHeading{Berufsspezifische rechtliche Grundlagen}  ~
  
  
  
  \begin{inparadesc}
    \item[Stundenanzahl:] 3
    \item[Lehrkraft:]  Jurist, fachkompetente Person
  \end{inparadesc}
  
  \begin{enumerate}\item
    
    Aufgaben und Kompetenzen des  Rettungssanitäters
    \\
    \noindent\FullRef{topic:kompetenzen-sanitaeter},
    \FullRef{topic:fortbildungspflicht}
    \item    Dokumentation im Rettungswesen (Einsatzprotokoll, Leitstellendokumentation, Transportnachweis)
    \\
    \noindent\FullRef{topic:dokumentation-rechtlich}
    \item    Hilfs- und Rettungswesen
    \\
    \FullRef{topic:vetrag-rettungswesen}
    \item    Straßenverkehrsordnung
    \\
    \noindent\FullRef{topic:strassenverkehrsordnung}
    \item    Patientenrechte
    \\
    integrativ in \FullRef{topic:langer-dienst}
    %\noindent\FullRef{topic:patientenrechte}
    \item    Grundlagen des Haftungsrechtes
    \\
    \FullRef{topic:vetrag-rettungswesen},
    \FullRef{topic:strafbarkeitspruefung},
    \FullRef{topic:schadenersatz}
    %\noindent\FullRef{topic:haftungsrecht}
    \item    Unterbringungsgesetz
    \\
    \noindent\FullRef{topic:unterbringung-jus},
    \FullRef{topic:unterbringungsgesetz},
    \FullRef{topic:unterbringung}
    \item    Reversfähigkeiten und Effekten
    \\
    \FullRef{topic:jus-fall-aufklaerung}
    %\FullRef{topic:effekte}
    \item    Mitnahme von Begleitpersonen
    \\
    wird extern abgehandelt
    %\noindent\FullRef{topic:begleitpersonen-mitnahme}
  \end{enumerate}
  
  \item \TabHeading{Anatomie und Physiologie}
  
  (Bestandteil des Sachgebiets 'Sanitätshilfe')
  
  \begin{inparadesc}
    \item[Stundenanzahl:] 4
    \item[Lehrkraft:]  Notarzt, Arzt für Allgemeinmedizin, approbierter Arzt,
    Facharzt, Turnusarzt,   Angehörige  des  gehobenen  Dienstes für Gesundheits-
    und Krankenpflege, Lehrsanitäter, fachkompetente Person
  \end{inparadesc}
  
  \begin{enumerate}
    \item    Blutkreislauf - Grundzüge
    \\
    \FullRef{topic:blutkreislauf}
    \item    Gliedmaßen - Grundzüge
    \\
    Obere Extremität:
    \FullRef{topic:schulterguertel},
    \FullRef{topic:obere-extremitaet},
    \\
    untere Extremität:
    \FullRef{topic:beckenguertel},
    \FullRef{topic:untere-extremitaet}
    \item    Haut -- Grundzüge
    \\
    \FullRef{topic:haut}
    \item    Schädel und Rumpf - Grundzüge Skelett
    \\
    Schädel: \FullRef{topic:schaedel}
    \item    Brustkorb - Grundzüge
    \\
    Thorax: \FullRef{topic:brustkorb-skelett}, \\
    Herz: \FullRef{topic:herz} \\
    Atemwege/Lunge: \FullRef{topic:atemwege-lunge}\\
    Speiseröhre: \FullRef{topic:speiseroehre}
    \item    Bauchraum - Grundzüge
    \\
    Gastro-intestinaltrakt:
    \FullRef{topic:gastrointestinaltrakt}, \\ 
    andere Bauchorgane:
    \FullRef{topic:sonstige-bauchorgane}.
  \end{enumerate}
  
  \item \TabHeading{Störungen der  Vitalfunktionen  und Regelkreise und zu setzende Maßnahmen}
  
  (Bestandteil des  Sachgebiets  'Sanitätshilfe')
  
  \begin{inparadesc}
    \item[Stundenanzahl:]  8
    \item[Lehrkraft:]  Notarzt, Arzt für Allgemein- medizin, approbierter Arzt,
    Facharzt,  Lehr- sanitäter,  fachkompe- tente Person
  \end{inparadesc}
  
  \begin{enumerate}\item
    
    Definition Vitalfunktion
    \\
    \noindent\FullRef{topic:vitalfunktionen-definition},
    \FullRef{topic:vitalwerte}
    \item    NACA-Schema
    \\
    \FullRef{topic:naca}
    \item    Bewusstsein
    \\
    (Bewusstseinsstörungen, Bewusstlosigkeit)
    \\
    \noindent Bewusstsein
    allgemein: \FullRef{topic:bewusstsein};
    Bewusstseinsstörungen:
    \FullRef{topic:bewusstseinsstoerung},
    \\
    Bewusstlosigkeit: \FullRef{topic:bewusstseinsstoerung}
    \item
    Atmung
    \begin{enumerate}
      \item      (Atemstörungen,
      \\
      \noindent\FullRef{topic:atemstörungen},
      \FullRef{topic:respiratorische-notfaelle}
      \item      Atemstillstand)
      \\
      \noindent\FullRef{topic:atemstillstand},
      \FullRef{topic:atemstillstand-aed}
    \end{enumerate}
    \item    Kreislauf
    \begin{enumerate}
      \item      (Kreislaufstörungen,
      \\
      Blutdruck: \FullRef{topic:blutdruck},\\
      Herz-Kreislaufstörungen:
      \FullRef{topic:herz-kreislauf-stoerungen},\\ 
      Schock:
      \FullRef{topic:schock-ehi}, \FullRef{topic:schock}
    \end{enumerate}
    \item    Kreislaufstillstand)
    \\
    \FullRef{topic:kreislaufstillstand},
    \FullRef{topic:kreislaufstillstand-aed}
    \item    Regelkreise
    \begin{enumerate}
      \item      Wärmehaushalt
      \\
      \FullRef{topic:temperaturregulation}
      \item      Wasser- und Elektrolythaushalt
      \\
      \FullRef{topic:wasser-elektrolythaushalt}
      \item      \Säure-Basen-Haushalt
      \\
      \FullRef{topic:saeure-basen-haushalt}
      \item      Stoffwechsel
      \\
      \FullRef{topic:stoffwechsel}
    \end{enumerate}
    \item
    Starke Blutung
    \\
    \FullRef{topic:starke-blutung},
    \FullRef{topic:druckverband},
    \FullRef{topic:schock-hypovolaemischer}
    \item
    Schock (Ursachen,
    \FullRef{topic:schock}
    \begin{enumerate}
      \item      Wirkung, Schockformen,
      \\
      \FullRef{topic:schock}, \FullRef{topic:schock-arten}
      \item      Verlauf, Schockzeichen)
      \\
      \FullRef{topic:schock-allgemeine-symptomatik}
    \end{enumerate}
    \item Akute Störung der Atmung
    \begin{enumerate}
      \item      Atembehinderung
      \\
      \noindent\FullRef{topic:atemstörungen}
      \item      Verlegung der Atemwege  durch Fremdkörper
      \\
      \noindent\FullRef{topic:atemstörungen},
      \FullRef{topic:mechanische-atemwegsverlegung}
      \item      Verlegung durch Schwellung
      \\
      \noindent\FullRef{topic:atemstörungen},
      \FullRef{topic:asthma-bronchiale}
      \item      Akute Atemnot
      \\
      \noindent\FullRef{topic:atemstörungen},
      \FullRef{topic:respiratorische-notfaelle}
      \item      Absaugung
      \\
      \FullRef{topic:absaugung}
      \item      Sauerstoff (Inhalationsrichtlinien, Dosierungsvorschriften)
      \\
      Sauerstoff: \FullRef{topic:sauerstoff},\\
      Beatmungshilfen: \FullRef{topic:beatmung-hilfsmittel}
      \item      Assistierte Beatmung (Indikationen, Durchführung)
      \\
      \FullRef{topic:beatmung-assistiert}
    \end{enumerate}
    \item    Feststellung des Todes
    \\
    \FullRef{topic:tod}
  \end{enumerate}\item \TabHeading{Notfälle bei verschiedenen Krankheitsbildern  und zu setzende Maßnahmen}
  
  (Bestandteil des Sachgebiets 'Sanitätshilfe')
  
  \begin{inparadesc}
    \item[Stundenanzahl:]  6
    \item[Lehrkraft:]  Notarzt, Arzt für Allgemeinmedizin, approbierter Arzt,
    Facharzt, Turnusarzt, Lehrsanitäter, fachkompetente Person
  \end{inparadesc}
  \begin{enumerate}\item
    Koma aus vorerst unbekannter Ursache
    \begin{enumerate}
      \item      (Schlaganfall,
      \\
      \FullRef{topic:insult}
      \item      Meningitis,
      \\
      \FullRef{topic:meningitis}
      \item      Diabetes,
      \\
      \FullRef{topic:diabetes-mellitus}
      \item      Vergiftung)
      \\
      \FullRef{topic:vergiftung}
    \end{enumerate}
    \item
    Krampfanfall \\
    \noindent\FullRef{topic:zerebrale-krampfanfaelle},
    \FullRef{topic:krampfanfaelle-kinder},
    \FullRef{topic:Wundstarrkrampf} \\
    \begin{enumerate}
      \item      (Epilepsie, Tetanie)
      \\
      \noindent\FullRef{topic:zerebrale-krampfanfaelle}
    \end{enumerate}
    \item Akuter Gefäßverschluss an den Gliedmaßen
    \begin{enumerate}
      \item      Venenthrombose,
      \\
      \FullRef{topic:venenthrombose}
      \item      Arterielle Embolie)
      \\
      \FullRef{topic:arterieller-gefaessverschluss}
    \end{enumerate}
    \item Pulmonale Notfälle
    \begin{enumerate}
      \item      Asthma bronchiale
      \\
      \FullRef{topic:asthma-bronchiale}
      \item      Lungenödem
      \\
      \FullRef{topic:lungenoedem}
      \item      Lungenembolie
      \\
      \FullRef{topic:lungenembolie}
      \item      Lungenentzündung
      \\
      \FullRef{topic:lungenentzuendung}
    \end{enumerate}
    \item Cardiale Notfälle
    \begin{enumerate}
      \item
      (Herz-Rhythmus-Störungen,
      \\
      \FullRef{topic:herzrhythmusstoerungen}
      \item
      Herzversagen,
      Linksherzschwäche,
      Rechtsherzschwäche,
      \\
      \FullRef{topic:herzinsuffizienz},
      \FullRef{topic:linksherzinsuffizienz},
      \FullRef{topic:rechtsherzinsuffizienz}
      \item
      Angina pectoris,
      \\
      \FullRef{topic:angina-pectoris}
      \item
      Herzinfarkt,
      \\
      \FullRef{topic:herzinfarkt}
      \item
      Hochdruckkrise)
      \\
      \FullRef{topic:hypertensive-krise}, \FullRef{topic:hypertensiver-notfall}
    \end{enumerate}
    \item Allgemeinchirurgische Notfälle
    \begin{enumerate}
      \item
      Akutes Abdomen
      \\
      \FullRef{topic:akutes-abdomen}
      \item
      Gastritis, Ulcus
      \\
      \FullRef{topic:gastritis}
      \item  Gastroenteritis,
      \\
      \FullRef{topic:gastroenteritis}
      \item  Pankreatitis,
      \\
      \FullRef{topic:pankreatitis}
      \item   Appendizitis,
      \\
      \FullRef{topic:appendizitis}
      \item   Gallenkolik,
      \\
      \FullRef{topic:gallenkolik}
      \item   Ileus,
      \\
      \FullRef{topic:ileus}
      \item    Mesenterialinfarkt,
      \\
      \FullRef{topic:mesenterialinfarkt}
      \item  Lebensmittelvergiftung,
      \\
      \FullRef{topic:lebensmittelvergiftung}
      \item gastrointestinale Blutung)
      \\
      \FullRef{topic:gastrointestinale-blutung},\\
      \FullRef{topic:gastrointestinale-blutung},
      \FullRef{topic:oesophagusvarizenblutung}
    \end{enumerate}
    \item Gynäkologische und urologische Erkrankungen
    \begin{enumerate}
      \item Nierenkolik,
      \\
      \FullRef{topic:nierenkolik}
      \item Nierenbeckenentzündung,
      \\
      \FullRef{topic:nierenbeckenentzuendung}
      \item Harnwegsinfekt,
      \\
      \FullRef{topic:harnwegsinfekt}
      \item akutes Harnverhalten,
      \\
      \FullRef{topic:harnverhalten}
      \item chronische Niereninsuffienz und Hämodialyse,
      \\
      Niereninsuffizienz:
      \FullRef{topic:niereninsuffizienz};
      
      Elektrolytstörungen durch
      Dialyse: \FullRef{topic:dialyse-elektrolyte};
      
      Dialyse als Therapie bei chronischer
      Niereninsuffizenz: \FullRef{topic:cni-therapie-dialyse};
      
      Dialysepatienten, besondere
      Gefahren: \FullRef{topic:dialysepatienten-gefahren}.
      \item extrauterine Gravidität,
      \\
      \FullRef{topic:extrauterine-graviditaet}
      \item Abortus,
      \\
      \FullRef{topic:abort}
      \item  Unterleibsblutungen,
      \\
      \FullRef{topic:unterleibsblutungen}
      \item Ovarialtumore,
      \\
      fraglich
      \item Vergewaltigung (typische Verletzungen, psychische Betreuung)
      \\
      wird extern abgehandelt
      %\FullRef{topic:vergewaltigung}
    \end{enumerate}
  \end{enumerate}
  
  \item \TabHeading{Spezielle Notfälle und zu setzende Maßnahmen  }
  
  (Bestandteil des Sachgebiets  'Sanitätshilfe')
  
  \begin{inparadesc}
    \item[Stundenanzahl:]  15
    \item[Lehrkraft:]   Notarzt, Arzt für  Allgemein-  medizin,  approbierter
    Arzt, Facharzt,   Turnusarzt,  Lehr-  sanitäter,  fachkompetente Person
  \end{inparadesc}
  
  \begin{enumerate}
    \item Traumatologische Notfälle
    \begin{enumerate}
      \item      Schädel-Hirn-Verletzungen,
      \\
      \FullRef{topic:sht}, \FullRef{topic:sht-diagnostik}
      \item      Hirnblutung
      \\
      \FullRef{topic:hirnblutung},
      \FullRef{topic:sht-diagnostik}
      \item      Hirndruck
      \\
      Hirndruckzeichen: \FullRef{topic:hirnblutung},
      \FullRef{topic:sht-diagnostik};
      
      Hirnstammeinklemmung bei
      Hirndrucksteigerung: \FullRef{topic:hirnblutung},
      \FullRef{topic:sht-diagnostik};
      
      Ursachen:
      \FullRef{topic:hirnblutung}, \FullRef{topic:insult};
      
      Lagerung: \FullRef{topic:hirnblutung}.
      \item
      Halswirbelsäulen- und  Wirbelsäulentrauma
      \\
      
      \begin{enumerate}
      	\item      (Sturzhelmabnahme,
      	\\
      	\FullRef{topic:sturzhelmabnahme}	
      	\item      HWS-Scheine Beschreibung (MP): \\
     	 Material \FullRef{topic:hws-schiene-beschreibung}, \\
      	Anwendung \FullRef{topic:hws-schiene-anlegen}
      
      	\item     
      	Motorik-Durchblutung-Sensibilitätskontrolle,
      	\\
      	\FullRef{topic:sehen-hoeren-fuehlen},
      	\FullRef{topic:ersteinschaetzung}
      	\item      Body-Check, Umgang mit  Schaufeltrage und
      	Vakuummatratze, Sandwich-Technik)\\
      	Body-Check: \FullRef{topic:traumacheck}, \\
      	Material:
      	\FullRef{topic:schaufeltrage-beschreibung},
      	\FullRef{topic:vakuummatratze-beschreibung}, 
      	\\ Anwendung:
      	\FullRef{topic:schaufeltrage-anwendung}, 
      	\FullRef{topic:vakuummatratze-anwendung}
      \end{enumerate}
      
      \item Extremitätentrauma
      \FullRef{topic:extrimitaetentrauma}
      \begin{enumerate}
        
        \item (Verletzungsarten,
        \FullRef{topic:extrimitaetentrauma}
        
        \item Motorik-Durchblutung-Sensibilitätskontrolle,\\
        \FullRef{topic:sehen-hoeren-fuehlen},
        \FullRef{topic:ersteinschaetzung}
        
        \item Stiefelgriff,
        Prinzip der Schienung,
        Ruhigstellung,\\
        \FullRef{topic:vacuumbeinschiene-beschreibung},
        \FullRef{topic:vakuummatratze-anwendung},
        \FullRef{topic:frakturzeichen},
        \FullRef{topic:knochenbruch-eh},
        \FullRef{topic:verstauchung-tra},
        \FullRef{topic:verstauchung-eh},
        \FullRef{topic:verrenkung-eh}
        \item        Schienung des Armes
        \\
        \FullRef{topic:schienung-arm}
        \item        Schienung des Beines
        \\
        Material:
        \FullRef{topic:vacuumbeinschiene-beschreibung},\\
        Anwendung: 
        \FullRef{topic:vacuumbeinschiene-anwendung}
        \item        Pneumatische Schiene
        \\
        wird extern abgehandelt
        \item        Vakuumschiene
        \\
        Vakuumbeinschiene,\\
        Material:
        \FullRef{topic:vacuumbeinschiene-beschreibung},\\
        Anwendung: 
        \FullRef{topic:vacuumbeinschiene-anwendung}
        \item        Extensionsschiene
        \\
        wird extern abgehandelt
      \end{enumerate}
      \item Thoraxtrauma
      \begin{enumerate}
        \item
        offene, geschlossene Brustkorbverletzung
        \\
        \FullRef{topic:thoraxtrauma}
        \item        Serienrippenbruch
        \\
        \FullRef{topic:serienrippenfraktur}
        \item        Pneumothorax
        \\
        \FullRef{topic:pneumothorax}
        \item        Spannungspneumothorax
        \\
        \FullRef{topic:pneumothorax}
      \end{enumerate}
      \item      Bauchtrauma
      \begin{enumerate}
        \item        offene, geschlossene Bauchverletzung
        \\
        \FullRef{topic:bauchtrauma-offen},
        \FullRef{topic:bauchtrauma-geschlossen}
        \item        Verletzung der Harn- und
        Geschlechtsorgane
        \\
        wird extern abgehandelt
        \item        Beckentrauma
        \\
        \FullRef{topic:beckentrauma}
      \end{enumerate}
      \item      Polytrauma (Definition, Prioritäten, Management)
      \\
      \FullRef{topic:polytrauma}
      \item      Wunden (mechanische, chemische, thermische)
      \\
      \FullRef{topic:wunden}
      \item      Dekubitus Prophylaxe und Lagerung bei Dekubitus
      \\
      \FullRef{topic:dekubitusprophylaxe}
      \item      Verbandlehre
      \\
      \FullRef{topic:druckverband},
      \FullRef{topic:starke-blutung}
      \item      Akut auftretende Blutungen
      \\
      \FullRef{topic:starke-blutung},
      \FullRef{topic:druckverband},
      \FullRef{topic:schock-hypovolaemischer}
      \item      Vergiftung (Ursachen und Verdacht, Aufnahmearten)
      \\
      \FullRef{topic:vergiftungen-ehi},
      \FullRef{topic:vergiftung}
    \end{enumerate}
    \item Psychiatrische Notfälle
    \begin{enumerate}
      \item      (Suizid,
      \\
      \FullRef{topic:suizid1}, \FullRef{topic:suizid2}
      \item      Psychose
      \\
      \FullRef{topic:psychose}
      \item      Suchterkrankungen und Entzugssyndrom
      \\
      \FullRef{topic:alkohol-entzug}
      %\FullRef{topic:sucht},
      %\FullRef{topic:entzugssysndrom}
      \item      Depression
      \\
      \FullRef{topic:depression}
      \item      Manie
      \\
      \FullRef{topic:manie}
    \end{enumerate}
    \item
    Schwangerschaft und Geburt
    \begin{enumerate}
      \item      Notfälle in der Schwangerschaft
      \\
      \FullRef{topic:notfall-fruehschwangerschaft}, 
      \FullRef{topic:notfall-spaetschwangerschaft}
      \item      Geburt
      \\
      \FullRef{topic:geburt}
      \item      Geburtskomplikationen
      \\
      \FullRef{topic:notfall-geburt}
      \item      Versorgung des Neugeborenen\\
      \FullRef{topic:geburt}
    \end{enumerate}
    \item Notfälle im Säuglings- und Kleinkindalter
    \begin{enumerate}
      \item      anatomische und physiologische Besonderheiten
      \\
      \FullRef{topic:kind-besonderheiten}
      \item      Kontrolle der Lebensfunktionen und lebensrettende Sofortmaßnahmen
      \\
      \FullRef{topic:pediatric-life-support}
      \item      Krampfanfälle
      \\
      Allgemein: \FullRef{topic:zerebrale-krampfanfaelle},\\
      Fieberkrampf: \FullRef{topic:krampfanfaelle-kinder},
      \\
      Wundstarrkrampf: \FullRef{topic:Wundstarrkrampf}
      \item      Pseudokrupp
      \\
      \FullRef{topic:laryngitis}
      \item      Epiglottitis
      \\
      \FullRef{topic:epiglottitis}
      \item      Keuchhusten
      \\
      \FullRef{topic:}
      \item      plötzlicher Kindstod
      \\
      \FullRef{topic:sids}
    \end{enumerate}
  \end{enumerate}
  \item \TabHeading{Defibrillation mit halbautomatischen Geräten}
  
  (Bestandteil des Sachgebiets  'Sanitätshilfe')
  
  \begin{inparadesc}
    \item[Stundenanzahl:]   8
    \item[Lehrkraft:]   Notarzt,  Arzt für Allgemeinmedizin,  approbierter Arzt,
    Facharzt, Lehrsanitäter, fachkompetente Person
  \end{inparadesc}
  
  \begin{enumerate}
    \item Der halbautomatische Defibrillator\\
    \FullRef{chp:AED}
    \item Handhabung eines halbautomatischen Defibrillators\\
    \FullRef{chp:AED}
    \item Gerätemanagement während der Reanimation\\
    \FullRef{chp:AED}
    \item Erfolgskontrolle\\
    \FullRef{chp:AED}
  \end{enumerate}
  \item \TabHeading{Gerätelehre und Sanitätstechnik }
  
  \begin{inparadesc}
    \item[Stundenanzahl:]  12
    \item[Lehrkraft:]   Lehrsanitäter, fachkompetente Person
  \end{inparadesc}
  
  \begin{enumerate}\item
    
    Medizinproduktegesetz -- Information
    \\
    \FullRef{topic:mpg-info}
    \item    Bergungs- und Lagerungstechniken
    \begin{enumerate}
      \item      Bergetuch
      \\
      \FullRef{topic:bergetuch}
      \item      Einheitskrankentrage
      \\
      \FullRef{topic:fahrtrage}
      \item      Tragsessel
      \\
      \FullRef{topic:tragsessel}
      \item      Fahrtrage
      \\
      \FullRef{topic:fahrtrage},
      \FullRef{topic:lagerungsarten}
      \item      Rollstuhl
      \\
      \FullRef{topic:rollstuhl}
    \end{enumerate}
    \item    Einsatzfahrzeug
    \\
    \FullRef{topic:einsatzmittel}
    \item    Beatmungsbeutel
    \\
    \FullRef{topic:beatmungsbeutel}
    \item    Absauggeräte
    \\
    \FullRef{topic:absaugung}
    \item    Sauerstoff
    \\
    \FullRef{topic:sauerstoff}
    \item    Infusionen und
    \\
    \FullRef{topic:infusion}, \FullRef{topic:venflon}
    \item    Infusionsgeräte
    \\
    \FullRef{topic:infusion}
    \item    Blutdruckmessung
    \\
    \FullRef{topic:rr-messung}
    \item    Stabilisierungs- und Schienungstechniken\\
    Material: \FullRef{topic:schienungsmaterial}, \\
    Anwendungen: \FullRef{topic:schienung}
    \begin{enumerate}
      \item      Stabilisierung der Halswirbelsäule
      \\
      Material \FullRef{topic:hws-schiene-beschreibung}, \\
      Anwendung \FullRef{topic:hws-schiene-anlegen}
      \item      Schaufeltrage
      \\
      Material: \FullRef{topic:schaufeltrage-beschreibung},
      \\ Anwendung: \FullRef{topic:schaufeltrage-anwendung}
      \item      Vakuummatratze
      \\
      Material: \FullRef{topic:vakuummatratze-beschreibung},
      \\ Anwendung:
      \FullRef{topic:vakuummatratze-anwendung}
    \end{enumerate}
    \item    Geburtenkoffer
    \\
    wird extern abgehandelt
    \item    Transportinkubator
    \\
    wird extern abgehandelt
  \end{enumerate}
  
  \item \TabHeading{Rettungswesen}
  
  \begin{inparadesc}
    \item[Stundenanzahl:]  4
    \item[Lehrkraft:]   Fachkompetente Person
  \end{inparadesc}
  
  \begin{enumerate}\item
    
    Rechtliche Grundlagen
    \\
    \FullRef{topic:rettungswesen-recht}
    \item    Zusammenarbeit mit  anderen Organisationen
    \\
    \FullRef{topic:rettungswesen-wien}
    \item    Einsatzarten
    \\
    \FullRef{topic:einsatzarten}
    \item    Transport- und Fahrzeugarten (Land,  Wasser, Luft)
    \\
    \FullRef{topic:einsatzmittel}
    \item    Normen, persönliche Schutzausrüstung
    \\
    wird extern abgehandelt
    \item    Fahrzeugausstattung (Dienstvorschriften)
    \\
    wird extern abgehandelt (TD/Fahrzeugeinschulung)
    \item    Rettungskette
    \\
    \FullRef{topic:rettungskette}
    \item    Hilfsfrist
    \\
    wird extern abgehandelt
    \item    Dienststellennetz
    \\
    \FullRef{topic:rettungswesen-wien}
    \item    Personal im Rettungsdienst
    \\
    \FullRef{chp:TAE}
    \item    Notarztsysteme
    \\
    \FullRef{topic:rettungswesen-wien},
    \FullRef{topic:einsatzmittel}
    \item    Leitstelle
    \\
    \FullRef{topic:leitstelle}
    \item    Kommunikationsarten
    \\
    \FullRef{topic:kommunikationswege}
    \item    Gefahren an der Einsatzstelle
    \\
    \FullRef{topic:gefahrenzonen}
    \item    Gefahrguteinsätze
    \\
    \FullRef{topic:gefahrengutunfall}
    \item    Sondertransporte
    \\
    wird extern abgehandelt
    
  \end{enumerate}
  \item \TabHeading{Katastrophen, Großschadensereignisse, Gefahrgutunfälle}
  
  \begin{inparadesc}
    \item[Stundenanzahl:]  4
    \item[Lehrkraft:]  Fachkompetente Person
  \end{inparadesc}
  
  \begin{enumerate}\item
    
    Katastrophen
    \begin{enumerate}
      \item
      Rechtliche Grundlagen,  Geltungsbereiche
      \\
      \FullRef{topic:katastrophe},
      \FullRef{topic:katastrophenhilfsgesetze}
      \item      Arten der Katastrophen
      \\
      \FullRef{topic:katastrophe},
      \FullRef{topic:katastrophenschutz}
      \item      Phasen der Katastrophenbewältigung
      \\
      \FullRef{topic:katastrophe-phasen}
      \item      Katastrophenhilfseinheiten
      \\
      \FullRef{topic:katastrophenschutz}
      \item      Führungsorganisation
      \\
      \FullRef{topic:katastrophe-fuehrung},
      \FullRef{topic:sanhist}
      \item      personelle, materielle und finanzielle Vorsorge
      \\
      \FullRef{topic:katastrophe-material}
      \item      Einsatzgrundsätze
      \\
      \FullRef{topic:einsatzgrundsaetze}
      \item      generelle Einsatzrichtlinien
      \\
      \FullRef{topic:einsatzrichtlinie}
      \item      Grundzüge der Triage
      \\
      \FullRef{topic:triage}
    \end{enumerate}
    \item
    Großschadensereignisse
    \begin{enumerate}
      \item
      Rechtliche Grundlagen
      \\
      \FullRef{topic:grossschaden}
      \item      Einstufung
      \\
      \FullRef{topic:grossschaden}
      \item      Alarmierung
      \\
      \FullRef{topic:alarmierung}
      \item      Schadensraum
      \\
      \FullRef{topic:sanhist}
      \item      Schadensplatz, Sicherheitseinrichtungen
      \\
      \FullRef{topic:sanhist}
      \item      Organisation beim Großunfall
      \\
      \FullRef{topic:sanhist}
      \item      Grundzüge der Triage
      \\
      \FullRef{topic:triage}
      \item      österreichisches Patientenleitsystem
      \\
      \FullRef{topic:pls}
      \item      Material und Ausrüstung
      \\
      \FullRef{topic:katastrophe-material}
      \item      Kommunikation
      \\
      \FullRef{topic:khd-organisation}
    \end{enumerate}
    \item
    Gefahrgutunfälle
    \begin{enumerate}
      \item      Arten von Gefahrgutunfällen
      \\
      \FullRef{topic:gefahrenzonen}
      \item      Gefahrzettel
      \\
      \FullRef{topic:gefahrzettel}
      \item      Gefahrsymbole
      \\
      \FullRef{topic:gefahrsymbole}
      \item      Warntafel
      \\
      \FullRef{topic:warntafel},
      \FullRef{topic:gefahrnummer}
      \item
      Verhalten am Unfallort
      \\
      \FullRef{topic:einsatzrichtlinie}
      \item      Koordination mit anderen Einsatzorganisationen
      \\
      wird extern abgehandelt
      \item      Absperrmaßnahmen
      \\
      \FullRef{topic:absperrmassnahmen}
      \item      Sofortmaßnahmen
      \\
      \FullRef{topic:einsatzrichtlinie}
    \end{enumerate}
  \end{enumerate}\item \TabHeading{Angewandte Psychologie und Stressbewältigung}
  
  
  
  \begin{inparadesc}
    \item[Stundenanzahl:]  4
    \item[Lehrkraft:]  Fachkompetente Person
  \end{inparadesc}
  
  \begin{enumerate}\item
    
    Belastung, Anforderung, Beanspruchung, workflow
    \\
    \FullRef{topic:belastung}
    \item    Überforderung, Unterforderung
    \\
    wird extern abgehandelt (PEER-Vortrag)
    \item    Beanspruchungsfolgen
    \\
    wird extern abgehandelt (PEER-Vortrag)
    \item    Stressursachen, -entstehung und -faktoren
    \\
    \FullRef{topic:stress}
    \item    Stressauswirkung
    \\
    \FullRef{topic:stress}
    \item    Früherkennung
    \\
    \FullRef{topic:stress}
    \item    Grundsätze der Stressvermeidung
    \\
    \FullRef{topic:stress}
    \item    Maßnahmen zur Verhütung und Verminderung von Beanspruchungsfolgen
    \\
    wird extern abgehandelt (PEER-Vortrag)
    \item    Psychische Betreuung von Kranken/Verletzten
    \\
    \FullRef{topic:kommunikationsregeln},
    \FullRef{topic:umgang-psy}
    \item    Gesprächsführung
    \\
    \FullRef{topic:kommunikationsregeln},
    \FullRef{topic:umgang-psy}
    \item    Vertrauensaufbau und Patienteninformation
    \\
    \FullRef{topic:kommunikationsregeln},
    \FullRef{topic:umgang-psy}
    \item    psychische Belastungssyndrome
    \\
    \FullRef{topic:ptss}
    \item    verwirrte Patienten
    \\
    wird integrativ in allen Kapitel abgehandelt
    \item    Begleitung und Betreuung Sterbender
    \\
    wird extern abgehandelt
    \item    Supervision
    \\
    wird extern abgehandelt (PEER-Vortrag)
  \end{enumerate}
  \item \TabHeading{Praktische Übungen ohne Patientenkontakt}
  
  \begin{inparadesc}
    \item[Stundenanzahl:] 16
    \item[Lehrkraft:] Lehrsanitäter, fachkompetente Person
  \end{inparadesc}
  
  \begin{enumerate}
    \item Regloser Notfallpatient
    \\
    \TakeNotes
    \item    Kontrolle der  Lebensfunktionen  (erwachsener Notfallpatient)
    \\
    \TakeNotes
    \item    Notfalldiagnose Bewusstlosigkeit (stabile Seitenlage)
    \\
    \TakeNotes
    \item    Notfalldiagnose Atemstillstand (Beatmung)
    \\
    \TakeNotes
    \item    Notfalldiagnose Kreislaufstillstand (Beatmung und Herzmassage)
    \\
    \TakeNotes
    \item    Blutstillung (Fingerdruck, Abdrückstellen, Druckverband, Abbindung, Amputatsversorgung)
    \\
    \TakeNotes
    \item    Blutdruckmessung
    \\
    \TakeNotes
    \item    Schockbekämpfung (Lagerungsarten)
    \\
    \TakeNotes
    \item
    Halswirbelsäulen- und Wirbelsäulentrauma
    \\
    \begin{enumerate}
      \item      (Sturzhelmabnahme,
      \\
      \TakeNotes
      \item      Halswirbelsäulenschienung,
      \\
      \TakeNotes
      \item      Motorik-Durchblutung-Sensibilitätskontrolle,
      \\
      \TakeNotes
      \item      Body-Check,
      \\
      \TakeNotes
      \item      Umgang mit Schaufeltrage,
      \\
      \TakeNotes
      \item      Umgang mit Vakuummatratze,
      \\
      \TakeNotes
      \item      Sandwich-Technik)
      \\
      \TakeNotes
    \end{enumerate}
    \item
    Extremitätentrauma
    \begin{enumerate}
      \item
      (Motorik-Durchblutung-Sensibilitätskontrolle,
      \\
      \TakeNotes
      \item      Stiefelgriff,
      \\
      \TakeNotes
      \item      Ruhigstellung,
      \\
      \TakeNotes
      \item      Schienung des Armes,
      \\
      \TakeNotes
      \item      Schienung des Beines,
      \\
      \TakeNotes
      \item      pneumatische Schiene,
      \\
      \TakeNotes
      \item      Vakuumschiene,
      \\
      \TakeNotes
      \item      Extensionsschiene)
      \\
      \TakeNotes
    \end{enumerate}
    \item
    Verbandlehre
    \\
    \TakeNotes
    \item Notfälle im Säuglings- und Kleinkindalter
    \begin{enumerate}
      \item
      (Kontrolle der Lebensfunktionen und lebensrettende Sofortmaßnahmen)
      \\
      \TakeNotes
    \end{enumerate}
    \item An- und Auskleiden, Körperpflege und Hygiene, Nahrungs- und Flüssigkeitsaufnahme, Harn- und Stuhlentleerung,
    Erbrechen
    \\
    \TakeNotes
    \item Ergonomische und schonende Arbeitsweise (Richtiges Heben und Tragen)
    \\
    \TakeNotes
    \item Handhabung der in Einsatzfahrzeugen zu verwendenden Geräte  (insbesondere
    Krankentrage, Tragsessel,  Sauerstoffgeräte, Kommunikationseinrichtungen) sowie
    die  Handhabung von Rollstühlen und Gehhilfen
    
    
  \end{enumerate}
\end{enumerate}

\end{multicols}

\end{WidePar}

\section{Berücksichtigte Lehrmeinungen}

\subsection{Arbeiter-Samariter-Bund Österreichs}

\subsubsection{Bundesverband}

\begin{WidePar}
\small
\begin{tabulary}{\linewidth}{LcLcL}
\toprule\addlinespace
\noindent\TabHeading{Kurztitel}
&
\noindent\TabHeading{Nr.}
&
\noindent\TabHeading{Referenz}
&
\noindent\TabHeading{seit}
&
\noindent\TabHeading{Querverweise}
\\\addlinespace\midrule\addlinespace
EH Reanimation%KT
&
01/2011%Nr
&
\bibentry{AsboeLehrmeinungEh2011-01} \cite{AsboeLehrmeinungEh2011-01}
%bibentry,cite
&
0.18.0%seit
&
\FullPrettyRef{topic:reanimation-eh}%FullPrettyRef
\\\addlinespace
EH Helmabnahme%KT
&
02/2011%Nr
&
\bibentry{AsboeLehrmeinungEh2011-02} \cite{AsboeLehrmeinungEh2011-02}%bibentry,cite
&
%seit
&
%FullPrettyRef
\\\addlinespace
RD Reanimation%KT
&
02/2011%Nr
&
\bibentry{AsboeLehrmeinungRd2011-02} \cite{AsboeLehrmeinungRd2011-02}%bibentry,cite
&
0.18.0%seit
&
\prettyref{chp:BLS}, \prettyref{chp:ALS}%FullPrettyRef
\\\addlinespace
EH Verbrennung%KT
&
03/2011%Nr
&
\bibentry{AsboeLehrmeinungEh2011-03} \cite{AsboeLehrmeinungEh2011-03}%bibentry,cite
&
%seit
&
%FullPrettyRef
\\\addlinespace
RD Nitrogabe%KT
&
03/2011%Nr
&
\bibentry{AsboeLehrmeinungRd2011-03} \cite{AsboeLehrmeinungRd2011-03}%bibentry,cite
&
%seit
&
%FullPrettyRef
\\\addlinespace
RD \ce{O2} bei MCI und ROSC%KT
&
04/2011%Nr
&
\bibentry{AsboeLehrmeinungRd2011-04} \cite{AsboeLehrmeinungRd2011-04}%bibentry,cite
&
%seit
&
%FullPrettyRef
\\\addlinespace
RD Not-/ schnelle Rettung%KT
&
05/2011%Nr
&
\bibentry{AsboeLehrmeinungRd2011-05} \cite{AsboeLehrmeinungRd2011-05}%bibentry,cite
&
%seit
&
%FullPrettyRef
\\\addlinespace
RD Rautekgriff%KT
&
06/2011%Nr
&
\bibentry{AsboeLehrmeinungRd2011-06} \cite{AsboeLehrmeinungRd2011-06}%bibentry,cite
&
0.18.0%seit
&
\FullPrettyRef{topic:rautekgriff}%FullPrettyRef
\\\addlinespace
Sturzhelmabnahme
&
07/2011
&
\bibentry{AsboeLehrmeinungRd2011-07} \cite{AsboeLehrmeinungRd2011-07,ASBOEAkademieUpdate2010Sturzhelm}
&
0.16.0
&

\\\addlinespace
RD \ce{O2} allgem.%KT
&
09/2011%Nr
&
\bibentry{AsboeLehrmeinungRd2011-09} \cite{AsboeLehrmeinungRd2011-09}%bibentry,cite
&
0.18.0%seit
&
\FullPrettyRef{m:sauerstoffberieselung}%FullPrettyRef
\\\addlinespace
Spineboard%KT
&
10/2011
&
\bibentry{AsboeLehrmeinungRd2011-10} \cite{AsboeLehrmeinungRd2011-10}%bibitem
&
 0.18.0
&
%FullPrettyRef
\\\addlinespace
RD \ce{O2}-Gabe Asthma%KT
&
11/2011%Nr
&
\bibentry{AsboeLehrmeinungRd2011-11} \cite{AsboeLehrmeinungRd2011-11}%bibentry,cite
&
0.18.0%seit
&
\FullPrettyRef{m:asthmaanfall}%FullPrettyRef
\\\addlinespace
Boa-Rettung%KT
&
14/2011
&
\bibentry{AsboeLehrmeinungRd2011-14} \cite{AsboeLehrmeinungRd2011-14}%bibitem
&
ausständig
&
%FullPrettyRef
\\\addlinespace
RD Guedel-Tubus%KT
&
15/2011%Nr
&
\bibentry{AsboeLehrmeinungRd2011-15} \cite{AsboeLehrmeinungRd2011-15}%bibentry,cite
&
%seit
&
%FullPrettyRef
\\\addlinespace
RD Larynxtubus%KT
&
16/2011%Nr
&
\bibentry{AsboeLehrmeinungRd2011-16} \cite{AsboeLehrmeinungRd2011-16}%bibentry,cite
&
%seit
&
%FullPrettyRef
\\\addlinespace
Absaugen Neugeborener%KT
&
18/2011
&
\bibentry{AsboeLehrmeinungRd2011-18} \cite{AsboeLehrmeinungRd2011-18}%bibitem
&
0.18.0%seit
&
\FullPrettyRef{m:neugeborenes}%FullPrettyRef
\\\addlinespace
Abnabelung%KT
&
19/2011%Nr
&
\bibentry{AsboeLehrmeinungRd2011-19} \cite{AsboeLehrmeinungRd2011-19}%bibitem,cite
&
%seit
&
%FullPrettyRef
\\\addlinespace
Geburt, Dammschutz%KT
&
20/2011%Nr
&
\bibentry{AsboeLehrmeinungRd2011-20} \cite{AsboeLehrmeinungRd2011-20}%bibentry,cite
&
0.18.0%seit
&
\ref{topic:dammschutz}, \ref{m:geburt}%FullPrettyRef
\\\addlinespace
Verbrennung
&
21/2011
&
\bibentry{AsboeLehrmeinungRd2011-21} \cite{AsboeLehrmeinungRd2011-21}

\bibentry{ASBOEAkademieUpdate2010Verbrennung} \cite{ASBOEAkademieUpdate2010Verbrennung}
&
0.16.0
&
\FullPrettyRef{m:verbrennung-eh}, \FullPrettyRef{m:verbrennung-leicht}, \FullPrettyRef{m:verbrennung-schwer}
\\\addlinespace
Notfallkompetenzen%KT
&
23/2011%Nr
&
\bibentry{AsboeLehrmeinungRd2011-23} \cite{AsboeLehrmeinungRd2011-23}%bibitem,cite 
&
in Arbeit%seit
&
\ref{m:hypertensiver-notfall}%FullPrettyRef
\\\addlinespace
%KT
&
%Nr
&
%bibentry,cite
&
%seit
&
%FullPrettyRef
\\\addlinespace
%KT
&
%Nr
&
%bibentry,cite
&
%seit
&
%FullPrettyRef
\\\addlinespace\bottomrule
\end{tabulary}


\end{WidePar}


\subsubsection{Landesverband Wien}

\subsubsection{Gruppe Floridsdorf-Donaustadt}
